%% Ejemplo de la plantilla de latex para las diapositivas del Máster en IA 

%%%%%%%%%%%%%%%%%%%%%%%%%%%%
%% Valores para el ratio de aspecto: 43, 169 
\documentclass[10pt, envcountsect, presentation, aspectratio=169]{beamer}
%%%%%%%%%%%%%%%%%%%%%%%%%%%%

%%%%%%%%%%%%%%%%%%%%%%%%%%%%
%% Incluir fichero de estilo del máster
\input plantillaMasterIA.sty
%%%%%%%%%%%%%%%%%%%%%%%%%%%%

%%%%%%%%%%%%%%%%%%%%%%%%%%%%
%% Información que aparecerá en la portada
\title[Nombre]{Razonamiento Aproximado Basado en Lógica Difusa}
\subtitle{Informe 2 de Conocimiento y Razonamiento Aproximado} % short title for footer

\author[E1. Apellidos1, E2. Apellidos2] % Para el pie de página, poner los autores abreviados separados por comas.
{%
	\sc{Estudiante1 Apellidos1}\\  % Autor 1
%	\textit{Ciencia de la Computación e Inteligencia Artificial}\\ % Area de conocimiento
%	\textit{Dept. Ingeniería de la Información y las Comunicaciones}\\  % Departamento
	\\
	\sc{Estudiante2 Apellidos 2}\\ % Autor 2
%	\textit{Ciencia de la Computación e Inteligencia Artificial}\\  % Area de conocimiento
%	\textit{Dept. Ingeniería de la Información y las Comunicaciones}\\  % Departamento
	\\   % No dejes espacio entre esta línea y la llave siguiente 
}

\institute[DIIC]% % Poner las iniciales de los diferentes departamentos de los autores separadas por comas.
{
	\textit{Universidad de Murcia}
}

\date{2025/2026} % Curso académico

%%%%%%%%%%%%%%%%%%%%%%%%%%%%%%%%%%
%% Contenido de las diapositivas
\setbeamerfont{normal text}{size=\normalsize} % Modifica el tamaño del texto de las diapositivas.
\AtBeginDocument{\usebeamerfont{normal text}}

%%%%%%%%%%%%%%%%%%%%%%%%%%%%%%%%%%

\begin{document}	

%%%%%%%%%%%%%%%%%%%%%%%%%%%%%%%%%%

%%%%%%%%%%%%%%%%%%%%%%%%%%%%%%%%%%

\begin{frame}[plain]
	\titlepage
\end{frame}

%%%%%%%%%%%%%%%%%%%%%%%%%%%%%%%%%%


\begin{frame}{Título del tema}{Título de la sección}
% \vskip -1cm % Para subir (-) o bajar el texto
\textbf{Título de la subsección (o del concepto en su caso)}
	\begin{itemize}
	\item Texto ...
	\item Texto ... \citep{Russell2020} ...
		\begin{itemize}
		\normalsize % Quítalo si quieres que se reduzca automáticamente el tamaño del texto	
		\item[--] Texto ...
		\item[--] La \href{https://es.wikipedia.org/wiki/Inteligencia_artificial}{Inteligencia Artificial} ...
			\begin{itemize}
			\normalsize % Quítalo si quieres que se reduzca automáticamente el tamaño del texto
			\item[$\bullet$] Texto ...
			\end{itemize}	
			\end{itemize} 
	\item Texto ...		
	\item Texto ...
	\end{itemize}
\end{frame}



%%%%%%%%%%%%%%%%%%%%%%%%%%%%%%%%%%
%%%%%%%%%%%%%%%%%%%%%%%%%%%%%%%%%%
% BIBLIOGRAFÍA
%%%%%%%%%%%%%%%%%%%%%%%%%%%%%%%%%%
%%%%%%%%%%%%%%%%%%%%%%%%%%%%%%%%%%

% Redefinimos la apariencia
\setbeamertemplate{headline}{}
\setbeamertemplate{frametitle}{\vskip 0.5cm	\small{\textbf{\insertframetitle}}}
\setbeamertemplate{footline}{}

% Mostramos la bibliografía
% \begin{frame}[allowframebreaks] %% Si ocupa más de una página
\begin{frame} %% Si hay sólo una página
	\frametitle{Referencias}
	\footnotesize
	% \vskip -5cm %% Para ajustar la distancia al comienzo
	\bibliographystyle{abbrvnat}
	\bibliography{references}	
\end{frame}    
		
%%%%%%%%%%%%%%%%%%%%%%%%%%%%%%%%%%
%%%%%%%%%%%%%%%%%%%%%%%%%%%%%%%%%%

\end{document}

%%%%%%%%%%%%%%%%%%%%%%%%%%%%%%%%%%
%%%%%%%%%%%%%%%%%%%%%%%%%%%%%%%%%%